\section{Introduction}

  Say that we have some abstract group $G$. Groups are defined purely by how elements multiply, which can make them hard to get a grip on. Linear algebra, on the other hand, is something we understand extremely well. So here is the idea of representation theory: \textit{study a group by letting it act linearly on a vector space}. Each group element becomes a matrix, group multiplication becomes matrix multiplication, and then we can bring in all of the tools of linear algebra (eigenvalues, traces, decompositions) to learn about $G$. 

  Therefore, the main questions we will ask are the following. 
  \begin{enumerate}
    \item What are all the ways that $G$ can act linearly on a vector space? 
    \item Can we break each such action into smaller, irreducible pieces, like prime factorization? 
    \item How can we tell two representations apart? (This will be answered with characters.) 
  \end{enumerate}

  Unless stated otherwise, all vector spaces are finite-dimensional over $\mathbb{C}$ and all groups are finite. Lie groups and Lie algebras are defined in my notes on differential geometry, and we will refer to them there. 
